% =====================================================================
%  ch0_front.tex — 封面、译者说明、目录
% =====================================================================
\thispagestyle{plain}
\begin{center}
{\LARGE\bfseries 三价图的拉格朗日 Floer 理论\\[0.3em]与曲线的同调镜像对称}\\[1.2em]
{\large Denis Auroux\ \ ·\ \ Alexander I. Efimov\ \ ·\ \ Ludmil Katzarkov}\\[0.8em]
{\normalsize 原载 \emph{Selecta Mathematica, New Series} 30 (2024), Paper 95}\\[0.2em]
{\normalsize arXiv:2107.01981v3 \quad 2024 年 7 月 5 日修订版}
\end{center}

\vspace{1.2em}
\begin{center}
{\large\bfseries 中文译注本}
\end{center}

\vspace{0.6em}
\noindent\rule{\textwidth}{0.4pt}
\subsection*{译者说明}
\addcontentsline{toc}{section}{译者说明}

本译本依据 arXiv:2107.01981v3 的 \LaTeX{} 源文件翻译，并回原文逐符号核对。原文共 59 页、6 幅图，图均由 Ti\emph{k}Z 代码给出，本译本保留其几何内容并重排。

\textbf{翻译原则}（依本库《数学文献翻译纪律》v2.1）：
\begin{enumerate}
\item 数学正确性第一；译文只对原文负责，不替作者补写、不替作者改错。原文确有笔误或明显跳步之处，一律保留原文表述，另以 \texttt{\% [Tutor Note: ...]} 批注说明（批注不进入 PDF）。
\item 术语先定后用：全书术语以附录《术语对照表》为准；有通行中文译名者用中文，无通行译名者保留英文，首次出现括注。
\item 公式逐符号核对，不用「大意如此」的转述代替公式；所有定义、定理、命题、引理、注的编号与原文一致（按节编号，共享计数）。
\item 引理、定理的证明不省略、不压缩；原文以「显然」「类似地」带过之处照译，必要处加批注。
\end{enumerate}

\textbf{符号约定}：$\C$ 复平面，$\PP^1$ 复射影直线，$\C^*=\C\setminus\{0\}$；$K$ 为 Novikov 域（系数在 $\R$ 中、指数有下界的 $T$ 的幂级数）；$\val$ 为 $K$ 上的赋值；$\F(M)$ 为构形 $M$ 的 Fukaya 范畴；$\Perf(X)$ 为 $X$ 上完美复形的导出范畴；$\Coh(X)$ 为凝聚层范畴。原文的 $A_\infty$ 符号约定（尤其是 $\mu^k$ 的符号与次数移位）在 §3.2 首次出现处加注说明。

\textbf{批注}：全书共 \textbf{11} 条译者批注，分「易错」「点拨」「直觉」三类，均带来源，以 \texttt{\% [Tutor Note]} 形式写在源码中，不进入排版结果。

\begin{flushright}
数译·SY05-27\\2026 年 11 月
\end{flushright}

\noindent\rule{\textwidth}{0.4pt}

\vspace{0.8em}
\newpage
\tableofcontents
\newpage
